% ============================================================================
\clearpage
\uxsection{Anexo A. Journey map en formato de tabla}\label{sec:anexo}

Se reproduce el contenido completo del mapa en forma tabular, de modo que sea consultable con
independencia de la resolución de la figura y accesible para lectura asistida.

\begin{uxtabla}{C{1.5cm}Z{1.0}Z{0.9}Z{1.1}}{Journey map de Laura Méndez, etapas 1 a 4}
\uxheadrow \thd{Etapa} & \thd{Acciones y puntos de contacto} & \thd{Pensamientos y emoción} & \thd{Punto de fricción y oportunidad} \\ \midrule
1. Disparo \newline \textit{Antes del uso} & Recibe la solicitud, identifica concepto y corte, decide dónde buscar. Mensajería, correo, calendario &
\emph{¿Es lo mismo del mes pasado? ¿A quién pregunto?} \newline Emoción: $-1$ &
El punto de partida es la memoria y los contactos, no el portal. \newline Enlace directo desde la mensajería y consultas guardadas visibles al abrir \\
2. Acceso \newline \textit{Durante el uso} & Inicia sesión; el sistema monta el espacio de su rol. Pantalla de acceso, espacio por rol &
\emph{Solo quiero llegar a la búsqueda.} \newline Emoción: $0$ &
Cualquier paso extra se paga caro antes de que la tarea empiece. \newline Sesión persistente que aterriza en el buscador \\
3. Búsqueda \newline \textit{Durante el uso} & Escribe el concepto en lenguaje de negocio, evalúa y descarta resultados. Buscador, tarjetas, filtros &
\emph{Hay dos que se llaman casi igual.} \newline Emoción: $-1$ a $+1$ &
Variables similares y miedo a tener que explicar el error. \newline Fuente, propietario y vigencia en la tarjeta; marca de versión retirada \\
4. Validación \newline \textit{Momento de la verdad} & Abre la ficha, lee definición, nota tribal y fecha. Catálogo, propietario, conocimiento tribal &
\emph{¿Ya lo revisó alguien o soy la primera?} \newline Emoción: $+2$ o $-2$ &
La definición vive en un correo o en la cabeza de alguien. \newline Ficha única con definición, versiones y notas del propietario \\
\end{uxtabla}

\begin{uxtabla}{C{1.5cm}Z{1.0}Z{0.9}Z{1.1}}{Journey map de Laura Méndez, etapas 5 a 7}
\uxheadrow \thd{Etapa} & \thd{Acciones y puntos de contacto} & \thd{Pensamientos y emoción} & \thd{Punto de fricción y oportunidad} \\ \midrule
5. Profundización \newline \textit{Durante el uso} & Filtra por periodo y consulta al asistente. Filtros, detalle, asistente con herramienta visible &
\emph{¿De dónde sacó el asistente ese número?} \newline Emoción: $+1$ &
Filtros técnicos en una consulta simple; respuestas sin fuente. \newline Revelación progresiva y cita de la fuente en cada cifra \\
6. Entrega \newline \textit{Durante el uso} & Copia el dato con referencia y responde. Copiado con procedencia, correo &
\emph{Quiero que se vea de dónde salió.} \newline Emoción: $+2$ &
Copiar pierde el contexto y anula la validación anterior. \newline Copiado con procedencia y enlace permanente al resultado \\
7. Reutilización \newline \textit{Después del uso} & Guarda la consulta; recibe el aviso de cambio y revisa el impacto. Consultas guardadas, notificación, historial &
\emph{¿El reporte que mandé sigue siendo válido?} \newline Emoción: $+1$ &
Nadie se entera del cambio hasta que dos áreas discrepan. \newline Aviso dirigido con alcance retroactivo declarado \\
\end{uxtabla}
